% !TEX program = xelatex
\documentclass[12pt, a4paper]{report}

% 字體與中文支援
\usepackage{xeCJK}
\setCJKmainfont{Microsoft JhengHei} % Windows 預設微軟正黑體，確保可編譯
\usepackage{geometry}
\geometry{a4paper, margin=2.54cm}

% 圖表與程式碼套件
\usepackage{graphicx}
\usepackage{listings}
\usepackage{xcolor}
\usepackage[hidelinks]{hyperref} % 隱藏目錄紅框
\usepackage{titlesec}
\usepackage{lipsum}

% 重新定義目錄與標題名稱為繁體中文
\renewcommand{\contentsname}{目錄}
\renewcommand{\listfigurename}{圖目錄}
\renewcommand{\listtablename}{表目錄}
\renewcommand{\figurename}{圖}
\renewcommand{\tablename}{表}
\renewcommand{\lstlistingname}{程式碼}
\renewcommand{\lstlistlistingname}{程式碼目錄}
\renewcommand{\chaptername}{第 \thechapter 章}

% 設定 Listings 程式碼排版格式 (具備左側行號)
\lstset{
    basicstyle=\ttfamily\small,
    numbers=left,
    numberstyle=\tiny\color{gray},
    stepnumber=1,
    numbersep=10pt,
    backgroundcolor=\color{white},
    showspaces=false,
    showstringspaces=false,
    showtabs=false,
    frame=single,
    rulecolor=\color{black},
    tabsize=4,
    captionpos=b,
    breaklines=true,
    breakatwhitespace=false,
    keywordstyle=\color{blue}\bfseries,
    commentstyle=\color{teal},
    stringstyle=\color{red}
}

\begin{document}

% ==================== 自訂期中考專屬封面頁與評分表 ====================
\begin{titlepage}
    \centering
    \vspace*{0.5cm}
    {\LARGE 專業排版與設計 期中考題 \par}
    \vspace{1.5cm}
    {\Huge \textbf{LINE Messaging API 教學文件製作} \par}
    \vspace{2cm}

    {\Large
    \renewcommand{\arraystretch}{1.5}
    \begin{tabular}{l}
    % ⚠️ 請把下方的「你的座號」改改成你實際的數字 ⚠️
    座號：\underline{ 68 \hspace{3.5cm} } \\   
    班級：\underline{ 資管3A \hspace{3.5cm} } \\
    學號：\underline{ 414637305 \hspace{2.8cm} } \\
    姓名：\underline{ 蔡竣宇 \hspace{3.7cm} } \\
    \end{tabular}
    \par}
    
    \vspace{1.5cm}
    
    % 加入考題要求的評分表
    \renewcommand{\arraystretch}{1.6}
    \begin{tabular}{|l|c|c|p{7.5cm}|}
    \hline
    評分項目 & 配分 & 得分 & 評分說明 \\ \hline
    \LaTeX 頁數達標 & 10\% & & 文件達五頁以上，內容充實、有條理 \\ \hline
    \LaTeX 排版美觀與結構 & 15\% & & 章節層次清晰，目錄可正常顯示 \\ \hline
    \LaTeX 進行 Github 版本控管 & 15\% & & 文件已成功上傳至 Github，並有良好的版本控制，至少有 3 個 effective commit \\ \hline
    \LaTeX 格式正確 & 20\% & & 文件可正常編譯 PDF，段落、圖片、程式碼 (需自動產生行號) 排版完善 \\ \hline
    \LaTeX 內容完整性 & 20\% & & 是否完整描述註冊過程與設定步驟，具技術深度 \\ \hline
    \LaTeX 程式碼品質與解說 & 20\% & & 程式碼可運作並有充分註解，並標註行號 \\ \hline
    \end{tabular}
    \vfill
\end{titlepage}

% 目錄頁
\pagenumbering{roman}
\tableofcontents
\newpage
\listoffigures
\newpage
\listoftables
\newpage
\lstlistoflistings
\newpage

\pagenumbering{arabic}

% ==================== 第 1 章 ====================
\chapter{LINE Messaging API 簡介與應用場景}
\section{技術簡介}
LINE Messaging API 是一套由 LINE 官方提供的開發工具介面，允許開發者建立客製化的聊天機器人（Chatbot）。透過這套 API，伺服器能夠雙向地與 LINE 使用者進行互動，實現自動化客服、即時通知、甚至互動式遊戲等功能。

\section{常見應用場景}
透過 webhook 接收使用者的訊息後，系統可以觸發多樣化的服務。為了說明其廣泛的應用，表 \ref{tab:usecases} 列出了目前業界常見的應用情境。

\begin{table}[h]
\centering
\caption{LINE Messaging API 應用場景整理}
\label{tab:usecases}
\begin{tabular}{|l|l|}
\hline
\textbf{領域} & \textbf{應用說明} \\ \hline
電子商務 & 訂單狀態查詢、推播專屬優惠券、購物車未結帳提醒 \\ \hline
客戶服務 & 24小時自動化常見問題解答 (FAQ)、轉接真人客服 \\ \hline
物聯網 (IoT) & 結合 ESP32 進行遠端家電控制與狀態回報 \\ \hline
\end{tabular}
\end{table}

\newpage
% ==================== 第 2 章 ====================
\chapter{LINE Developers 帳號註冊申請與 Channel 建立}
\section{帳號註冊}
開發者需要先前往 LINE Developers 官方網站註冊開發者帳號。可以使用一般個人的 LINE 帳號直接綁定登入，完成基本資料的填寫並同意服務條款後，即可進入開發者後台 (Console)。

\section{建立 Provider 與 Channel}
進入 Console 後，首先需要建立一個 Provider（提供者，通常為公司或專案名稱）。接著在 Provider 底下建立一個新的 Channel，請選擇「Messaging API」類型。

\begin{figure}[h]
\centering
% 使用黑色方塊代替截圖，確保沒有圖檔也能編譯成功
\rule{10cm}{5.5cm} 
\caption{LINE Developers Channel 建立示意畫面}
\label{fig:line_console}
\end{figure}

設定過程中需要填寫 Channel icon、Channel name 以及 Channel description。填寫完畢後點擊 Create 即可完成建立流程。

\newpage
% ==================== 第 3 章 ====================
\chapter{Channel Secret 及 Access Token 取得與設定}
\section{取得憑證資訊}
為了確保伺服器與 LINE 平台之間的通訊安全，我們需要兩把關鍵的鑰匙：
\begin{itemize}
    \item \textbf{Channel Secret}：位於 Basic settings 頁籤下方，主要用於驗證接收到的 Webhook 請求是否真的來自 LINE 官方。
    \item \textbf{Channel Access Token}：位於 Messaging API 頁籤下方，點擊 Issue 按鈕即可生成。我們的伺服器在主動發送或回覆訊息給使用者時，必須在 HTTP Header 中附上此 Token 作為授權證明。
\end{itemize}

\section{Webhook 的設定與測試}
為了讓本機端開發的程式能夠接收來自 LINE 的外部請求，我們可以使用 \texttt{ngrok} 將本機的 Port (例如 8000) 映射到具有 SSL 憑證的公開網址。
取得 ngrok 的 HTTPS 網址後，將其加上 \texttt{/callback} 路徑，填入 LINE Console 的 Webhook URL 欄位，並點擊 Verify 進行測試。

\newpage
% ==================== 第 4 章 ====================
\chapter{範例程式說明}
\section{LINE Bot 範例程式碼說明}
本章節展示如何使用 Python 與 Flask 框架實現一個簡單的回聲機器人 (Echo Bot)。此程式碼使用了官方的 \texttt{line-bot-sdk}，能夠接收使用者傳送的文字，並原封不動地回傳給使用者。

\subsection{app.py 解析}
這段程式碼會在 Port 8000 啟動一個網頁伺服器，並負責驗證 \texttt{X-Line-Signature}，確認簽章無誤後，交由 \texttt{handle\_message} 函式透過 \texttt{ReplyMessageRequest} 進行回覆。

% 匯入外部的 app.py 並加入程式碼目錄
\lstinputlisting[language=Python, caption={Python app.py 範例}, label={lst:app_py}]{app.py}

\newpage
% ==================== 第 5 章 ====================
\chapter{結論}
透過本文件的引導，我們完成了從 LINE Developers 帳號的註冊、Channel 的建立、重要憑證的獲取，一直到透過 Python 與 Flask 建立 Webhook 伺服器的完整流程。這為未來進階應用，例如結合資料庫、圖文選單 (Rich Menu) 或是 Flex Message 奠定了穩固的基礎。

\newpage
% ==================== 第 6 章 ====================
\chapter{本文件 LaTeX 原始碼}
依照考題要求，本章節將這份 LaTeX 文件的原始碼匯入，展示排版所使用的各項指令與套件配置。

% 匯入自己的 main.tex 並加入程式碼目錄
\lstinputlisting[language=TeX, caption={main.tex 原始碼}]{main.tex}

\end{document}